\documentclass{situationist-preprint}

\title{The Impossibility of Nothing}
\msid{MET-01-d}
\msfamily{MET-01}
\msdiscipline{Metaphysics}
\mstype{Consequence}
\mskeywords{nothingness; empty world; serious possibility; modal measure; stability}

\begin{document}
\maketitle

\begin{abstract}
MET-01-a argues that the empty world is possible in the weak sense that its description is consistent, but that it receives measure zero among the possibilities and is not a stable possibility. We draw the consequence for the claim that there might have been nothing. The claim is ambiguous between combinatorial possibility and serious possibility, a notion we define by measure. On the first reading it is true; on the second it is false. The question why there is something rather than nothing contrasts the actual world with an alternative that is not seriously possible, and in that sense nothing is impossible.
\end{abstract}

\section{Two readings}

Let $W_n$ be the space of assignments of existence and non-existence to $n$ potential concrete objects, with the measure that gives each assignment an equal share (MET-01-a, Definition~4.1 and P2), as in the reconstruction of the argument of \citet{vanInwagen1996} given there. Say that a set $S$ of worlds is \emph{negligible} if, for every $k \geq 1$, its measure in $W_n$ is less than $1/k$ for all sufficiently large $n$; and that a proposition is \emph{seriously possible} if the set of worlds at which it holds is not negligible.

\begin{proposition}\label{p:two}
The proposition that there is nothing concrete is combinatorially possible and not seriously possible.
\end{proposition}

\begin{proof}
By P1 of MET-01-a, every combination of existence and non-existence of the potential objects is possible, the empty combination included; and by Proposition~4.3 of MET-01-a the empty world belongs to every $W_n$. By Theorem~4.4 of MET-01-a, for every $k \geq 1$ and every $n \geq k$ the measure of the empty world is less than $1/k$, so the set consisting of the empty world is negligible.
\end{proof}

Serious possibility is not an invention for the occasion. It is the notion at work when a probabilist says that a fair coin tossed forever will almost surely not land heads every time: the all-heads sequence is a consistent outcome, and it is not a serious one. Proposition~\ref{p:two} says that the empty world stands to the space of possible worlds as the all-heads sequence stands to the space of infinite sequences. The claim that there might have been nothing is true if \enquote{might} expresses combinatorial possibility, and false if it expresses serious possibility.

\section{Which reading the question uses}

The question why there is something rather than nothing is contrastive. It asks why the actual world falls on one side of a partition rather than the other, and a contrastive question presupposes that the foil was a live alternative: we do not ask why the coin landed in some sequence or other rather than heads forever. On the combinatorial reading, the foil is possible, and the question looks as if it demands a reason that tipped reality away from it. On the serious reading, the foil is not live, and the demand lapses: there is nothing to be tipped away from. This is the sense in which the answer of MET-01 dissolves the contrast without denying the consistency of the empty world.

Proposition~\ref{p:two} also distinguishes the present impossibility from two stronger ones. \citet{Lowe1996} argues that the empty world is metaphysically impossible, because there are necessarily existing abstract objects whose existence depends on there being concrete ones. If REL-01 is correct, a necessary being exists and the empty world is impossible outright, as MET-01-a notes. Neither claim is needed here, and neither is entailed: the subtraction argument of \citet{Baldwin1996} may succeed in showing that the empty world is possible, and Proposition~\ref{p:two} stands. What is impossible, in the sense established, is not the consistency of nothing but its being a serious alternative.

\section{Instability}

The second consequence concerns stability. MET-01-a observes that every one of the $n$ neighbors of the empty world in $W_n$, the worlds that differ from it with respect to a single object, is non-empty, whereas a non-empty world has at least $n - 1$ non-empty neighbors when $n \geq 2$. A possibility is \emph{stable} if some small variation of it preserves its status. Every variation of the empty world is a variation into something. Combined with Proposition~\ref{p:two}, this yields the full sense of the title: nothing is a consistent combination that occupies a vanishing share of the space of possibilities and is surrounded only by alternatives in which something exists. It is possible in the way an isolated point of measure zero is possible, and in no stronger way.

\begin{thebibliography}{9}\small
\bibitem[Baldwin(1996)]{Baldwin1996} Baldwin, T. (1996). There might be nothing. \emph{Analysis}, 56(4), 231--238.
\bibitem[Lowe(1996)]{Lowe1996} Lowe, E.~J. (1996). Why is there anything at all? \emph{Proceedings of the Aristotelian Society, Supplementary Volume}, 70, 111--120.
\bibitem[van Inwagen(1996)]{vanInwagen1996} van Inwagen, P. (1996). Why is there anything at all? \emph{Proceedings of the Aristotelian Society, Supplementary Volume}, 70, 95--110.
\end{thebibliography}

\end{document}
